% =====================================================================
% 8. Conclusion
% =====================================================================
\section{Conclusion}
\label{sec:conclusion}

Holding default probabilities fixed, the $99.9\%$ VaR of the portfolio ranges from USD~\val{q8.M2.VaR_ref|mn2} million under Gaussian factors to USD~\val{q8.M1.VaR_ref|mn2} million under the empirical distribution, with the $t$-copula model only \val{q8.M3.VaR_vs_M2|pct1}\% above the Gaussian one. The tail of the loss is sensitive to the distribution of the market factors, but to their margins far more than to their copula: allowing the two indices to crash together matters little as long as each is Gaussian and the firm-specific shocks are Gaussian and independent. A natural extension would combine a $t$-copula with fat-tailed margins, obtaining the tail thickness of M1 without tying the $99.9\%$ quantile to two historical weeks.

The same message recurs elsewhere. Through time, the parametric models are stable because the thresholds are recalibrated, while the empirical model inherits the ghost effects of historical simulation. Against the regulatory benchmark, the model-based unexpected loss is \val{q10.ratio_UL_M2|f1} times the IRB capital, a matter of correlation levels, not of granularity or tail shape. The storm overlay, which acts on recoveries and is independent of the credit cycle by assumption, raises the expected loss by \val{q11.dEL_exact|ppm} parts per million and ES by about USD~\val{q11.first_order.M2.addon_tail_usd|usd2}, four orders of magnitude below the Monte Carlo error of the estimates. Concentration is irrelevant under independence and costs less than $0.02\%$ of ES even under dependence; the PD channel of \citet{Pozdyshev25} and a $20\%$ gust stress raise the effect several- and tenfold, yet ES moves by less than $0.01\%$.

\paragraph{Caveats.}
The empirical model depends on a handful of observations. With $10{,}000$ simulations the Monte Carlo error at $99.9\%$ is of the same order as some of the differences of interest, which is why our conclusions rest on two million scenarios. The comparison with the IRB formula is driven by the systematic weights $\lambda_k$ supplied with the synthetic portfolio; with weights in the regulatory range, model-based and regulatory capital would be much closer. Finally, the climate overlay covers one acute physical hazard under the present climate and through the collateral only: transition and chronic risks, liquidity and operational channels, and scenario-conditional losses in the sense of supervisory climate stress tests are outside its scope.
